\labeledsection{Decision Networks}{sec:dn}
We now know how to update our world model, represented as a \linkterm{set}{def:set} of \linkterm{random variables}{def:random_variable}, given observations. Now we need to talk about \textbf{taking actions}.



\begin{tcolorbox}[
    enhanced,
    attach boxed title to top left={xshift=5mm, yshift=-3mm, yshifttext=-1mm},
    colback=blue!5!white,
    colframe=blue!75!black,
    colbacktitle=red!80!black,
    title=Question,
    fonttitle=\bfseries,
    boxed title style={size=small, colframe=red!50!black} 
]
Given a world model and a \linkterm{set}{def:set} of actions, what are the likely consequences of action?  How good are these consequences? 
\end{tcolorbox}

\ideabox{
(1) Represent actions as \textit{special} \linkterm{random variables}{def:random_variable}.
\begin{itemize}
    \item Given disjoint actions $a_1, \cdots, a_n$, we introduce a \linkterm{random variable}{def:random_variable} $A$ with domain $\set{a_1, \cdots, a_n}$. Then we can query $\probd{X \mid A=a_i}$
\end{itemize}
(2) Assign numerical values to the outcomes of actions, i.e. a \linkterm{function}{function} $\func{u}{\text{dom}(X)}{\mathbb{R}}$

(3) Choose the action that \textbf{maximizes} the \textbf{expected value} if $u$
}


\anondefinition{
\defemph{Decision Theory} investigates \defemph{decision problems}, i.e. how a \linkterm{utility-based agent}{utility_based_agent} $a$ deals with choosing among actions based on the desirability of their outcomes given by a real-valued \defemph{utility function} $U$ on states $s \in S$, i.e. $\func{U}{S}{\mathbb{R}}$
}{def:decision_theory}

\commandnote{
If our states are \linkterm{random variables}{def:random_variable}, then we obtain a \linkterm{random variable}{def:random_variable} for the \linkterm{utility function}{utility_based_agent}. 

Let $\func{X_i}{\Omega}{D_i}$ \linkterm{random variables}{def:random_variable} on a \linkterm{probability space}{def:prob_space} $\tuple{\Omega, \mathcal{F}, P}$ and $\func{f}{\cartprod{D_1, \cdots, D_n}}{\mathbb{R}}$. Then $F(x):=f(X_0(x), \cdots, X_n(x))$ is a \linkterm{random variable}{def:random_variable} $\Omega \to \mathbb{R}$. Meaning, any \linkterm{function}{function} on the domains of \linkterm{random variables}{def:random_variable} is itself a \linkterm{random variable}{def:random_variable}.
}

\Example{
Consider commuting to work as en example experiment where we have $\Omega = \set{\text{accident on the road}, \text{minor crash on the road}, \text{smooth commute}}$. 

We define a \linkterm{random variable}{def:random_variable} $X$ representing travel time in minutes:
\begin{itemize}
    \item $X(\text{accident on the road})=90$
    \item $X(\text{minor crash on the road})=45$
    \item $X(\text{smooth commute})=20$
\end{itemize}
These are the states our world can be in. We want to define a utility function $F$ representing our happiness. Because $X$ is a \textbf{random}, $F$ becomes \textbf{random} too, we do not get one happiness level, we get a distribution of possible hapiness levels based on the underlying reality $\Omega$.

Say that $f(t) = 100 - t$ (the longer the trip, the less happy I am). Then we have: 

\begin{tabular}{|l|c|c|} \hline
Reality $(\omega)$ & Travel Time $(X(\omega))$ & Happiness $(f(X(\omega)))$ \\ \hline   
accident on the road & 90 & 10 \\ \hline
minor crash on the road & 45 & 55 \\ \hline
smooth commute & 20 & 80 \\ \hline  
\end{tabular}

}{}

\anondefinition{
Given a \linkterm{probability space}{def:prob_space} $\tuple{\Omega, \mathcal{F}, P}$ and a \linkterm{random variable}{def:random_variable} $X:\Omega \to D$ with $D \subseteq \mathbb{R}$, then 
$$\mathbb{E}(X):=\sum_{x \in D}P(X=x) \cdot x$$
is called the \defemph{Expected Value} of $X$. Analogously, let $e_1, \cdots, e_n$ be a sequence of \linkterm{events}{def:event}. Then the \defemph{Expected Value} of $X$ given $e_1,\cdots, e_n$ is defined as 
$$\mathbb{E}(X \mid e_1, \cdots, e_n) := \sum_{x \in D }P(X=x \mid e_1, \cdots, e_n) \cdot x$$
}{expected_value}

\Example{
Continuing with the Commute Time example, assume we have the following probabilities:
\begin{itemize}
    \item $P(X=20) = 0.7$
    \item $P(X=45) = 0.2$
    \item $P(X=90) = 0.10$
\end{itemize}

\begin{align*}
\mathbb{E}(X)&:=\sum_{x \in D}P(X=x) \cdot x \\
&:=P(X=20) \cdot 20 + P(X=45) \cdot 45 + P(X=90) \cdot 90 \\ 
&:=0.7 \cdot 20 + 0.2 \cdot 45 + 0.10 \cdot 90 \\ 
\mathbb{E}(X)&:=32 \text{ minutes}
\end{align*}
\begin{tcolorbox}[colback=red!5!white,colframe=red!75!black]
This means, (while I might never actually spend exatly $32$ minutes on any single day), if I commute for, e.g. a year, my average travel time will be $32$ minutes.
\end{tcolorbox}
}{}

\anondefinition{
Given:
\begin{itemize}
    \item $D= \set{a_1, \cdots, a_n}$: a set of possible action, where we have a special \linkterm{random variable}{def:random_variable} $A:\Omega \to D$ representing the action.
    \item $X_i:\Omega \to D_i$: \linkterm{random variables}{def:random_variable} (the states)
    \item $U:\cartprod{D, D_i} \to \mathbb{R}$: a \linkterm{utility function}{utility_based_agent}
\end{itemize}
Then, the \defemph{Expected Utility} of the action $a \in D$ is the \linkterm{expected value}{expected_value} of $U$ (interpreted as a \linkterm{random variable}{def:random_variable}) given $A=a$:
\[
\text{EU}(a) := \sum_{\tuple{x_1, \cdots, x_n} \in \cartprod{D_1, \cdots, D_n}} P(X_1 = x_1, \cdots, X_n=x_n \mid A=a) \cdot U(x_1, \cdots, x_n, a)
\]
}{expected_utility}

\Example{
\begin{itemize}
\item $D = \set{\text{take}, \text{leave}}$ for taking or leaving the umbrella.
\item $X$: random variable for the state of the weather outside with $D_x = \set{\text{rain}, \text{sum}}$ with \linkterm{priors}{def:bayesian_terms}: $P(X=\text{rain}) = 0.3, P(X=\text{sun})=0.7$
\item $U:\cartprod{D, D_x} \to R$ the utility function defined as:
\begin{itemize}
    \item $U(\text{sun}, \text{leave}) = 100$
    \item $U(\text{sun}, \text{take}) = 80$
    \item $U(\text{rain}, \text{leave}) = 0$
    \item $U(\text{rain}, \text{take}) = 60$
\end{itemize}
\end{itemize}

\begin{align*}
\text{EU}(\text{take}) &:= \sum_{x_i \in \set{\text{rain}, \text{sun}}} P(X = x_i \mid A=\text{take}) \cdot U(x_i,\text{take}) \\  
&:= P(X = \text{rain} \mid A=\text{take}) \cdot U(\text{rain},\text{take}) + P(X = \text{sun} \mid A=\text{take}) \cdot U(\text{sun},\text{take}) \\ 
&:= P(X = \text{rain}) \cdot U(\text{rain},\text{take}) + P(X = \text{sun}) \cdot U(\text{sun},\text{take}) \\ 
&:= 0.3 \cdot 60 + 0.7 \cdot 80 = 74 \\ 
\end{align*}

\begin{align*}
\text{EU}(\text{leave}) &:= \sum_{x_i \in \set{\text{rain}, \text{sun}}} P(X = x_i \mid A=\text{leave}) \cdot U(x_i,\text{leave}) \\  
&:= P(X = \text{rain} \mid A=\text{leave}) \cdot U(\text{rain},\text{leave}) + P(X = \text{sun} \mid A=\text{leave}) \cdot U(\text{sun},\text{leave}) \\ 
&:= P(X = \text{rain}) \cdot U(\text{rain},\text{leave}) + P(X = \text{sun}) \cdot U(\text{sun},\text{leave}) \\ 
&:= 0.3 \cdot 0 + 0.7 \cdot 100 = 70 \\ 
\end{align*}
}{}

\definition{MEU Principle for Rationality}{
We call an action \defemph{rational} if it \defemph{Maximizes the Expected Utility (MEU)}. A \linkterm{utility-based agent}{utility_based_agent} is called \linkterm{rational}{def:rational}, iff it always chooses a \linkterm{rational}{def:rational} action.
}{def:MEU}

\definition{Decision Network}{
A \defemph{decision network} is a \linkterm{Bayesian Network}{def:bayes_net} with two additional kinds of nodes:
\begin{enumerate}
    \item \defemph{Action Nodes}: representing a \linkterm{set}{def:set} of possible actions \hfill (square)
    \item A single \defemph{Utility Node} (also called value node) \hfill (diamond)
\end{enumerate}
Simple Example (no dependencies):

\begin{tikzpicture}[
    node distance=2cm,
    % Chance Node: Circle
    chance/.style={circle, draw=blue!75, fill=blue!5, very thick, minimum size=1.2cm},
    % Action Node: Square
    action/.style={rectangle, draw=green!40!black, fill=green!5, very thick, minimum size=1.2cm},
    % Utility Node: Diamond
    utility/.style={diamond, draw=red!75!black, fill=red!5, very thick, minimum size=1.2cm}
]

    % Place the nodes
    \node[chance] (W) {Weather};
    \node[action, below of=W] (A) {Umbrella};
    \node[utility, right of=A, xshift=1.5cm, yshift=1cm] (U) {U};

    % Draw the arrows
    % Weather affects Utility
    \draw[-Stealth, thick] (W) -- (U);
    % Action affects Utility
    \draw[-Stealth, thick] (A) -- (U);

    % Note: If the Action depended on observing the weather (e.g. looking out the window),
    % there would be a dashed arrow from Weather to Action.

\end{tikzpicture}
}{def:dn}

\begin{tcolorbox}[
    colback=blue!3!white,
    colframe=blue!70!white,
    title=General Algorithm]
Given \linkterm{evidence}{def:bayesian_terms} $E_j = e_j$ and \linkterm{action nodes}{def:dn} $A_1, \cdots, A_k$, compute the \linkterm{expected utility}{expected_utility} of each action, given \linkterm{evidence}{def:bayesian_terms}, for each action we have: 
\[
\text{EU}(a\mid E=e_j) := \sum_{\tuple{x_1, \cdots, x_n} \in \cartprod{D_1, \cdots, D_n}} 
P(X_1 = x_1, \cdots, X_n=x_n \mid A=a, E_j = e_j) \cdot 
U(x_1, \cdots, x_n, a)
\]

And we are interested in the single action that \linkterm{MEU}{def:MEU}:
\[
A_{\text{MEU}} = \argmax{a_1, \cdots, a_k} \text{EU}(a_i\mid E=e_j)
\]

And we say $\text{MEU} = \text{EU}(A_{\text{MEU}} \mid E=e_j)$ (maximum expected utility is the expected utility we get if we would choose the action that maximizes the expected utility)
\end{tcolorbox}

\commandnote{
\linkterm{Decision Theorists}{def:decision_theory} are mostly \textit{economists}, not statisticians. Economists usually talk about \textit{prizes} (which we call states).
}

\anondefinition{
Given prizes $A$ and $B$, an agent can express preferences of the form:
\begin{itemize}
    \item $A \succ B$: $A$ is preferred over $B$
    \item $A \sim B$: indifference between $A$ and $B$
    \item $A \succcurlyeq B$: $B$ is not preferred over $A$ (the agent wants $A$ at least as much as $B \leadsto$ \textit{weak preference})
\end{itemize}
}{preferences}

\commandnote{
In computer science, we define \linkterm{relations}{relation_on} $\succ, \sim$ on our \linkterm{set}{def:set} of states $\mathcal{S}$
}

\begin{tcolorbox}[colback=red!5!white,colframe=red!75!black]
But that works only for \linkterm{deterministic environments}{env_types}. In a \linkterm{non-deterministic}{env_types} one, the agent simply does not know the consequences of the actions. Now we look at what \linkterm{Decision Theorists}{def:decision_theory} call \textit{lotteries}
\end{tcolorbox}

\anondefinition{
Let $\mathcal{S}$ be a \linkterm{set}{def:set} of states. We call a \linkterm{random variable}{def:random_variable} $X$ with domain $\set{A_1, \cdots, A_n} \subseteq \mathcal{S}$ a \defemph{lottery} and write $[p_1, A_1 ; p_2, A_2 ; \cdots ; p_n, A_n]$ where $p_i = P(X=A_i)$
}{lottery}


\begin{tcolorbox}[
    colback=blue!3!white,
    colframe=blue!70!white,
    title=Idea]
A \linkterm{lottery}{lottery} represents the result of a \linkterm{nondeterministic}{env_types} action that can have outcomes $A_i$ with \linkterm{prior}{def:bayesian_terms} $p_i$. For the binary case, we use $[p,A ; 1-p, B]$. We can then extend \linkterm{preferences}{preferences} to include \linkterm{lotteries}{lottery}, as a measure of how \textit{strongly} we prefer one prize over another.

\begin{tikzpicture}[node distance=2cm, thick]

    % Define the Lottery node
    \node (L) at (0,0) {\Large $L$};

    % Define the Outcome nodes
    \node (A) at (3,1) {\Large $A$};
    \node (B) at (3,-1) {\Large $B$};

    % Draw the branches with probabilities
    \draw[-{Stealth}] (L.east) -- (A.west) node[midway, above] {\Large $p$};
    \draw[-{Stealth}] (L.east) -- (B.west) node[midway, below] {\Large $1-p$};

\end{tikzpicture}
\end{tcolorbox}

\commandnote{
We assume the \linkterm{set}{def:set} of states $\mathcal{S}$ to be closed under \linkterm{lotteries}{lottery}, ,i.e. \linkterm{lotteries}{lottery} are also states. That allows us to consider \linkterm{lotteries}{lottery} such as $[p,A ; 1-p,[q,B ; 1-q,C]]$

\begin{tikzpicture}[
    node distance=2.5cm, 
    thick,
    outcome/.style={circle, draw=black, fill=gray!10, minimum size=8mm}
]

    % Root Lottery
    \node (L1) at (0,0) {\Large $L_1$};

    % First level outcomes
    \node[outcome] (A) at (3,1.5) {$A$};
    \node (L2) at (3,-1.5) {\Large $L_2$}; % Nested lottery node

    % Second level outcomes (from L2)
    \node[outcome] (B) at (6,-0.5) {$B$};
    \node[outcome] (C) at (6,-2.5) {$C$};

    % Edges for first lottery [p, A; 1-p, L2]
    \draw[-{Stealth}] (L1) -- (A) node[midway, above left] {$p$};
    \draw[-{Stealth}] (L1) -- (L2) node[midway, below left] {$1-p$};

    % Edges for nested lottery [q, B; 1-q, C]
    \draw[-{Stealth}] (L2) -- (B) node[midway, above] {$q$};
    \draw[-{Stealth}] (L2) -- (C) node[midway, below] {$1-q$};
\end{tikzpicture}
}

\begin{tcolorbox}[
    colback=blue!3!white,
    colframe=blue!70!white,
    title=Preference Relations]
\begin{itemize}
    \item $X \succcurlyeq Y \land Y \succcurlyeq X \Rightarrow X \sim Y$
    \item $X \succcurlyeq Y \land \neg (Y \succcurlyeq X) \Rightarrow X \succ Y$
\end{itemize}
\end{tcolorbox}

\anondefinition{
We call a set $\succ$ of \linkterm{preferences}{preferences} \defemph{rational}, iff the following constraints hold:
\begin{itemize}
\item \defemph{Orderability}: $A \succ B \lor B \succ A \lor A \sim B$
\item \defemph{Transitivity}: $A \succ B \land B \succ C \Rightarrow A \succ C$
\item \defemph{Continuity}: $A \succ B \succ C \Rightarrow (\exists p.[p,A ; 1-p,C] \sim B)$
\item \defemph{Substitutability}: $A \sim B \Rightarrow [p,A ; 1-p,C] \sim [p,B;1-p,C]$
\item \defemph{Monotonicity}: $A \succ B \Rightarrow ((p > q) \Leftrightarrow [p,A;1-p,B] \succ [q,A;1-q,B])$
\item \defemph{Decomposability}: $[p,A;1-p,[q,B;1-q,C]] \sim [p,A;((1-p)q),B;((1-p)(1-q)),C]$
\end{itemize}
}{rational_preferences}

\begin{tcolorbox}[colback=red!5!white,colframe=red!75!black]
From a \linkterm{set}{def:set} of \linkterm{rational preferences}{rational_preferences}, we can obtain a meaningful utility function.
\end{tcolorbox}

In plain English, this definition says when a person's preferences are considered \emph{rational} in decision theory.
\begin{itemize}
    \item \defemph{Orderability}: for any two options, you must be able to compare them. You either prefer one, prefer the other, or are indifferent.
    \item \defemph{Transitivity}: if you prefer $A$ to $B$, and $B$ to $C$, then you should also prefer $A$ to $C$.
    \item \defemph{Continuity}: if $A$ is better than $B$, and $B$ is better than $C$, then there is some probability $p$ such that you are indifferent between $B$ for sure and the lottery $[p,A ; 1-p,C]$.
    \item \defemph{Substitutability}: if you are indifferent between $A$ and $B$, then replacing $A$ by $B$ inside any lottery should not change your preference.
    \item \defemph{Monotonicity}: if $A$ is better than $B$, then a lottery with a higher probability of $A$ should be preferred to one with a lower probability of $A$.
    \item \defemph{Decomposability}: a nested lottery should be equivalent to a single flat lottery with the same overall probabilities.
\end{itemize}
\begin{tikzpicture}[
    thick,
    outcome/.style={circle, draw=black, fill=gray!10, minimum size=8mm}
]

% Left: nested lottery [p,A ; 1-p,[q,B ; 1-q,C]]
\begin{scope}[xshift=0cm]
    \node (L1) at (0,0) {\Large $L_1$};
    \node[outcome] (A1) at (3,1.5) {$A$};
    \node (L2) at (3,-1.5) {\Large $L_2$};
    \node[outcome] (B1) at (6,-0.5) {$B$};
    \node[outcome] (C1) at (6,-2.5) {$C$};

    \draw[-{Stealth}] (L1) -- (A1) node[midway, above left] {$p$};
    \draw[-{Stealth}] (L1) -- (L2) node[midway, below left] {$1-p$};
    \draw[-{Stealth}] (L2) -- (B1) node[midway, above] {$q$};
    \draw[-{Stealth}] (L2) -- (C1) node[midway, below] {$1-q$};

    \node at (3,-3.6) {\small $[p,A ; 1-p,[q,B ; 1-q,C]]$};
\end{scope}

% Middle equivalence symbol
\node at (8.6,-1.1) {\Large $\sim$};

% Right: decomposed lottery [p,A ; (1-p)q,B ; (1-p)(1-q),C]
\begin{scope}[xshift=11cm]
    \node (L) at (0,0) {\Large $L$};
    \node[outcome] (A2) at (3,1.5) {$A$};
    \node[outcome] (B2) at (3,0) {$B$};
    \node[outcome] (C2) at (3,-1.5) {$C$};

    \draw[-{Stealth}] (L.east) -- (A2.west) node[midway, above] {$p$};
    \draw[-{Stealth}] (L.east) -- (B2.west) node[midway, above] {$(1-p)q$};
    \draw[-{Stealth}] (L.east) -- (C2.west) node[midway, below] {$(1-p)(1-q)$};

    \node at (1.5,-3.6) {\small $[p,A ; (1-p)q,B ; (1-p)(1-q),C]$};
\end{scope}

\end{tikzpicture}


\Theorem{
\textbf{Ramsey's Theorem}: Given \linkterm{rational preferences}{rational_preferences}, there exists a real valued function $U$ such that $U(A) \geq U(B)$, iff $A \succcurlyeq B$ and $U([p,S_1;\cdots;p_n, S_n]) = \sum_{i} p_i U(S_i)$
}{Ramsey_thm}

\commandnote{
This theorem states that if an agent's / person's \linkterm{preferences}{preferences} are \linkterm{rational}{rational_preferences}, then we can assign numbers to outcomes in a way that matches those \linkterm{preferences}{preferences}. i.e. there exists a utility \linkterm{function}{function} $U$ that gives each outcome a number. If $A \succcurlyeq B$ then $U(A) \geq U(B)$.  The higher the utility number, the more preferred the outcome. For a \linkterm{lottery}{lottery} $[p,S_1;\cdots;p_n, S_n]$, its utility is just the probability-weighted average of the utilities of its outcomes.
}


\anondefinition{
We call a \linkterm{linear (total) ordering}{linear_order} on states a \defemph{value function} or \defemph{ordinal utility function}. With \linkterm{deterministic}{env_types} prizes only (no \linkterm{lottery}{lottery} choices), only a \linkterm{total ordering}{linear_order} on prizes can be determined.
}{}

\commandnote{
If we don't need to care about \textit{relative} utilities of states, e.g. to compute non-trivial \linkterm{expected utilities}{expected_utility}, that's all we need anyway! i.e. when there are no probabilities to compare, “utility” is just a way to sort outcomes from best to worst, not a precise measurement.
}

\definition{
Standard approach to assessment of human utilities
}{
Compare a given state $A$ to a \defemph{standard lottery} $L_p$ that has:
\begin{itemize}
    \item best possible prize $u_{\top}$ with probability $p$
    \item worst possible catastrophe $u_{\bot}$ with probability $1-p$
\end{itemize}
adjust \linkterm{lottery}{lottery} probability $p$ until $A \sim L_p$. Then $U(A) = p$
}{standard_lottery}

\Example{
Choose $u_{\top} \hat{=}$current state, $u_{\bot} \hat{=}$ instant death

\begin{tikzpicture}[>=Stealth, thick]
    % The starting state/lottery label
    \node (L) at (0,0) {pay \$30 $\sim L$};
    % The outcome nodes
    \node[anchor=west] (A) at (4,0.6) {continue as before};
    \node[anchor=west] (B) at (4,-0.6) {instant death};
    % The branches
    % Adjusting the start point slightly to the right of the text
    \draw[->] (2.1, 0.15) -- (A.west) node[midway, above left, xshift=2mm] {0.999999};
    \draw[->] (2.1,-0.15) -- (B.west) node[midway, below left, xshift=2mm] {0.000001};
\end{tikzpicture}
}{}


\begin{tcolorbox}[
    colback=blue!3!white,
    colframe=blue!70!white,
    title=Popular Utility Functions]

\anondefinition{
\defemph{Normalized utilities}: $u_{\top}=1, u_{\bot}=0$ (not meaningful)
}{normalized_utilities}

\anondefinition{
\defemph{Money}: very intuitive, often easy to determine, but actually not well-suited as utility function
}{money_utility}

\anondefinition{
\defemph{Micromorts}: one millionth chance of instant death. (useful for paying to reduce product risk, etc.). But not necessarily a good measure of risk, if the risk is only severe injury or illness
}{micromorts}

\anondefinition{
\defemph{QALYs: quality adjusted life years} $\leadsto$ very better, more informative, useful for medical decisions involving substantial risk 
}{QALYs}
\end{tcolorbox}

\definition{Value of Perfect Information (VPI)}{
Let $A$ be the \linkterm{set}{def:set} of available actions and $F$ is a \linkterm{random variable}{def:random_variable}. Given \linkterm{evidence}{def:bayesian_terms} $E_i=e_i$, let $\alpha$ be the action that \linkterm{MEU}{def:MEU} a priori, and $\alpha_f$ the action that \linkterm{MEU}{def:MEU} given $F=f$, i.e.
\begin{align*}
\alpha &= \argmax{a \in A} \text{EU}(a \mid E_i = e_i) \\ 
\alpha_f &= \argmax{a \in A} \text{EU}(a \mid E_i = e_i , F=f)
\end{align*}
The \defemph{value of perfect information (VPI)} on $F$ given \linkterm{evidence}{def:bayesian_terms} $E_i = e_i$ is defined as:
\[
\text{VPI}_{E_i=e_i}(F) := \left( \sum_{f \in \text{dom}(F)} P(F=f \mid E_i = e_i) \cdot \text{EU}(\alpha_f \mid E_i=e_i, F=f) \right) - \text{EU}(\alpha \mid E_i=e_i)
\]
}{vpi}

\Example{
consider two blocks $A,B$ and exactly one has oil which worth $\mathdollar K$. You can drill in one location. You know the \linkterm{priors}{def:bayesian_terms}: $P(\text{oil at A})=P(\text{oil at B}) = 0.5$. 

We can model this a \linkterm{Decision Network}{def:dn} with one chance node (oil location), one \linkterm{action node}{def:dn} (drilling location), and one \linkterm{utility node}{def:dn}:

\begin{tikzpicture}[
    node distance=1.5cm and 2.5cm,
    thick,
    % Action Node: Rectangle
    action/.style={rectangle, draw, minimum width=2.5cm, minimum height=0.8cm},
    % Chance Node: Oval
    chance/.style={ellipse, draw, minimum width=2.5cm, minimum height=0.8cm},
    % Utility Node: Diamond
    utility/.style={diamond, draw, minimum size=1cm},
    % Table style
    table/.style={font=\small}
]

    % --- Nodes ---
    \node[action] (D) {DrillLoc};
    \node[chance, below=of D] (O) {OilLoc};
    \node[utility, right=of D, yshift=-0.75cm] (U) {U};

    % --- Edges ---
    \draw[-Stealth] (D) -- (U);
    \draw[-Stealth] (O) -- (U);

    % --- Utility Table ---
    \node[right=0.5cm of U, draw, inner sep=5pt] (UTable) {
        \begin{tabular}{c|c|c}
            $D$ & $O$ & $U(D, O)$ \\ \hline
            a & a & $K\$$ \\
            a & b & 0 \\
            b & b & $K\$$ \\
            b & a & 0
        \end{tabular}
    };

    % --- Probability Table ---
    \node[below left=0.1cm and 0.5cm of O] (PTable) {
        \begin{tabular}{c|c}
              & $P(O)$ \\ \hline
            a & 0.5 \\
            b & 0.5
        \end{tabular}
    };

\end{tikzpicture}

\begin{tcolorbox}[
    enhanced,
    attach boxed title to top left={xshift=5mm, yshift=-3mm, yshifttext=-1mm},
    colback=blue!5!white,
    colframe=blue!75!black,
    colbacktitle=red!80!black,
    title=Question,
    fonttitle=\bfseries,
    boxed title style={size=small, colframe=red!50!black} 
]
Without any information, what is the action that \linkterm{MEU}{def:MEU}?
\end{tcolorbox}
\begin{align*}
\text{EU}(D=a) &= P(O=a) \cdot U(D=a,O=a) + P(O=b) \cdot U(D=a,O=b) \\
&= 0.5 \cdot K + 0.5 \cdot 0 = \frac{K}{2} \\
\text{EU}(D=b) &= P(O=a) \cdot U(D=b,O=a) + P(O=b) \cdot U(D=b,O=b) \\
&= 0.5 \cdot 0 + 0.5 \cdot K = \frac{K}{2}
\end{align*}

\begin{tcolorbox}[
    enhanced,
    attach boxed title to top left={xshift=5mm, yshift=-3mm, yshifttext=-1mm},
    colback=green!5!white,      
    colframe=green!40!black,    
    colbacktitle=red!80!black,  
    title=Answer,
    fonttitle=\bfseries,
    boxed title style={size=small, colframe=red!50!black}
]
Which means, without any additional information, both actions have the same expected utility. We write $\text{MEU}(\emptyset)$ for the maximum expected utility given no extra information:

\[
\text{MEU}(\emptyset) = \argmax{a} \text{EU}(a) = \frac{K}{2}
\]
\end{tcolorbox}
\begin{tcolorbox}[
    enhanced,
    attach boxed title to top left={xshift=5mm, yshift=-3mm, yshifttext=-1mm},
    colback=blue!5!white,
    colframe=blue!75!black,
    colbacktitle=red!80!black,
    title=Question,
    fonttitle=\bfseries,
    boxed title style={size=small, colframe=red!50!black} 
]
What if someone offered to give us information about where the oil is. What is the value of that information? i.e. How much we are willing to pay to get such information?
\end{tcolorbox}
\begin{tcolorbox}[
    enhanced,
    attach boxed title to top left={xshift=5mm, yshift=-3mm, yshifttext=-1mm},
    colback=green!5!white,      
    colframe=green!40!black,    
    colbacktitle=red!80!black,  
    title=Answer,
    fonttitle=\bfseries,
    boxed title style={size=small, colframe=red!50!black}
]
Assume we have information that says \textit{the oil location is b} (our \linkterm{evidence}{def:bayesian_terms}), then we get EU of 0 for $D=a$ and we get EU of $K$ for $D=b$ and we will choose to dig in $b$ immediately.

The value of this perfect information is simply what we get with information ($K$) minus what we would get without it ($\frac{K}{2}$).

\[
\text{VPI}(E) = \text{MEU}(A \mid E) - \text{MEU}(\emptyset)
\]
\end{tcolorbox}
}{}
